%Hilfsmakros (Werden von Hauptmakros in Macros.tex verwendet):


%TODO-Gestaltungsmakro f�r Textmodus mit einer Unterscheidung, ob etwas �bergeben worden ist oder nicht
\newcommand{\myTodo}[1]{%
		\ifx\relax#1\relax% Falls Argument #1 leer:
			\textcolor{red}{\textbf{ToDo!}}%
		\else% Falls #1 mit Inhalt:
			\textcolor{red}{\textbf{ToDo:} #1}%
		\fi%
}

%TODO-Gestaltungsmakro f�r Mathe-Modus mit einer Unterscheidung, ob etwas �bergeben worden ist oder nicht
\newcommand{\myMathTodo}[1]{%
		\ifx\relax#1\relax% Falls Argument #1 leer:
			\textcolor{red}{\mathbf{\textnormal{ ToDo! }}}%
		\else% Falls #1 mit Inhalt:
			\textcolor{red}{#1}%
			%\textcolor{red}{\mathbf{\textnormal{ ToDo: }}{#1}}%
		\fi%
}



%[#1]: optional: Pfad zu einem bereits vorhandenem Bild, das �berarbeitet werden soll (in eckigen Klammern)
%{#2}: TODO-Beschreibung
\newcommand{\myImgTodo}[2][]
{
 % new figure floating environment
 \begin{figure}[!htb]
  \centering
  % TODO-image
		\ifx\relax#1\relax%
			\begin{tikzpicture}
				\draw[
					use as bounding box,
					black,
					thick,
					double, 
					double distance = 4pt, 
					rounded corners = 6pt] 
					(-4,-3)rectangle(4,3);
					\node[rotate=45,scale = 5,opacity = 0.5] at (0,0) {TODO!};
			\end{tikzpicture}
		\else%
			\includegraphics[width = 0.8\textwidth]{#1}\relax
		\fi
	\caption{\myTodo{#2}}
 \end{figure}
}
